\scanpage{222}
linéaire sur $\mathcal M$, on trouve une unique antidérivation $\delta$ de $\Lambda^*\mathcal M$ qui sur $\Lambda^1\mathcal M$ se réduise à $\varphi$, et on a $\delta^2=0$. (N.B. On a $H_0(\Lambda^*\mathcal M,\delta)=S/\varphi(\mathcal M)$, et si $\mathcal M$ est libre de t.f., alors $H_i(\Lambda^*\mathcal M,\delta)=0$ pour $i\ne0$ sss $\mathcal M\longrightarrow S$ est ``transversal'' -- il suffit p.ce. qu'une base \unsure{au moins|donnée} de $\mathcal M$ définisse une suite \uline{$S$-régulière}, et c'est aussi néc. si $S$ noeth et \deleted{\unsure{le}} $\varphi(\mathcal M)\subset\operatorname{rad} S$)
Ceci dit, on obtient les complexes (2) et (4) en prenant respectivement $S=\operatorname{Sym}^* M$ ou $S=\Gamma^*M$, $\mathcal M=M\otimes_A S$, et en prenant l'hom $\mathcal M\to S$ \deleted{\unsure{la}} dont la restriction à $M$ est l'hom can $M\xrightarrow{\sim}\deleted{\operatorname{Sym}^1M\subset S}\inserted{S^1\subset S}$ \deleted{\unsure{resp. $\Gamma^1M$}} ($S^1=\operatorname{Sym}^1 M$ resp $\Gamma^1 M$).

Les opérateurs $d\delta,\delta d$ dans $\operatorname{Sym}^*\otimes\Lambda^*$ resp $\Gamma^*\otimes\Lambda^*$ ne sont pas bien jolis, mais on voit de suite (tout \unsure{anticrochet|anticommutateur} des antidérivations étant une dérivation) que
\[
 \theta=d\delta+\delta d
\]
est une dérivation, \deleted{et} \deleted{qui dans le cas $S^{\unsure{1|*}}=\Gamma^*$ est compatible aux p.d.} ne changeant \deleted{\unsure{par} le} le bidegré, et qui sur \deleted{\unsure{$S^1$}} $S^1=\operatorname{Sym}^1$ ou $\Gamma^1$ et sur $\Lambda^1$ est l'identité; \inserted{de plus si $S^*=\Gamma^*$, il est compatible aux puiss. div. sur $\Gamma^*$.} Donc on en conclut facilement que a lieu
\[
 \boxed{(d\delta+\delta d)x=\theta x=nx}\qquad
 \substack{x\text{ de degré }\underline{\text{total}}\text{ }n\\ \text{i.e. }x\in\displaystyle\sum_{\substack{p+q=n\\\unsure{p,q\geq0}}}S^p\otimes\Lambda^q}
\]
Cela montre donc que si $A$ est de car. zéro (donc $n$ inversible $\forall\,n>0$) alors $H^{n}_{\mathrm{DR}}(S^*\otimes\Lambda^*)$ et $H^{\mathrm{Kosz}}_{n}(S^*\otimes\Lambda^*)$ ($S^*$ étant $\operatorname{Sym}^*\simeq\Gamma^*$) est nul si \deleted{\unsure{$d$}} $n>0$ (sans hyp. de platitude!). Plus gén., $H^n_{\mathrm{DR}}$ et $H_n^{\mathrm{Kosz}}$ sont nuls pour tout $n$ tel que $n.1_A$ inv. dans $A$.
\margin{$\Vert$}

% archive sheet number: 221

